% Capítulo "Direito ao Brincar" do relatório -- EDITÁVEL À MÃO (specs/2026-11-05_versao_final_site R5).
% Exportado uma vez por gera_latex.py --exportar-capitulos (2026-11-05); desde então o TEXTO é escrito aqui, não
% vem mais de relatorio/textos_curados.json (que é só do site). Os DADOS continuam gerados: as figuras
% (\figura{...}{../../visualizacoes/a4/...} e mapas/a4/) saem do analise.py; as tabelas no corpo
% (\input{gerado/tabelas_corpo/...}) e os rótulos do apêndice (\ref{tab:...}) saem do gera_latex.py, que confere
% a cada execução se tudo o que este arquivo cita existe.

\chapter{Direito ao Brincar}\label{cap:eixo-5}

\begin{achados}\begin{itemize}

\item O eixo é representado hoje só pela violência territorial, um indicador de contexto da população geral (todas as idades), de 2024

\item A taxa de homicídios varia de 2,8 por 100 mil habitantes (Rocinha) a 76,3 (Centro); sem o Centro, o maior valor é o de Madureira (33,5)

\item Os homicídios por ação policial são mais altos na Cidade de Deus (19,6 por 100 mil), na Penha (15,0) e em Irajá (12,9)

\item Os homicídios de jovens negros se concentram em Madureira (9,2 por 100 mil), Pavuna (6,8) e Cidade de Deus (6,5)

\end{itemize}\end{achados}

O direito ao brincar depende de espaços seguros para a convivência das crianças. Nesta edição, o eixo é tratado a partir da violência territorial, que afeta o uso das ruas, praças e equipamentos públicos pelas famílias. Os dados são do Instituto de Segurança Pública, por Região Administrativa, para 2024, e se referem à população de todas as idades: medem o contexto em que as crianças vivem, não a violência sofrida por elas.

\section{Violência territorial}

\figura{mapa}{../../mapas/a4/mapa_violencia_territorial_homicidios_ra_2024.pdf}{Taxa de homicídios por RA (2024) — população geral}{Elaboração IPP com dados de \citeonline{isp_violencia2024}. Limites de bairros: \citeonline{ipp_limites_bairros}. Sistema de referência SIRGAS 2000.}{mapa:mapa-violencia-territorial-homicidios-ra-2024}

Em 2024, a taxa de homicídios apresentou forte desigualdade entre as Regiões Administrativas do Rio de Janeiro, variando de 2,8 (Rocinha) a 76,3 homicídios por 100 mil habitantes (Centro); sem o Centro, o maior valor é o de Madureira (33,5). As maiores taxas aparecem concentradas em algumas áreas da porção central e norte do município, enquanto diversas regiões do litoral e da Zona Sul apresentam valores mais baixos. O indicador evidencia diferenças importantes no nível de violência territorial dentro da cidade.

Como o dado considera a população geral, de todas as idades, ele deve ser utilizado como medida do contexto de violência do território, e não como indicador direto de homicídios na primeira infância. No relatório, esse mapa integra o eixo de Direito ao Brincar, na seção de violência territorial.

\figura{mapa}{../../mapas/a4/mapa_violencia_territorial_homicidios_acao_policial_ra_2024.pdf}{Homicídios por ação policial por RA (2024) — população geral}{Elaboração IPP com dados de \citeonline{isp_violencia2024}. Limites de bairros: \citeonline{ipp_limites_bairros}. Sistema de referência SIRGAS 2000.}{mapa:mapa-violencia-territorial-homicidios-acao-policial-ra-2024}

Em 2024, a taxa de homicídios por ação policial apresentou forte variação entre as Regiões Administrativas do Rio de Janeiro. Os maiores valores foram registrados na Cidade de Deus (19,62 por 100 mil habitantes), Penha (14,97) e Irajá (12,93). Em contraste, algumas regiões apresentaram taxa igual a zero, como Copacabana, Rocinha, Complexo do Alemão e Santa Teresa.

O mapa evidencia, portanto, uma distribuição territorial bastante desigual desse tipo de violência. Assim como no indicador anterior, os dados se referem à população geral, de todas as idades, devendo ser interpretados como medida do contexto de violência territorial, e não como indicador específico da primeira infância.

\figura{mapa}{../../mapas/a4/mapa_violencia_territorial_homicidios_jovens_negros_ra_2024.pdf}{Homicídios de jovens negros por RA (2024) — população geral}{Elaboração IPP com dados de \citeonline{isp_violencia2024}. Limites de bairros: \citeonline{ipp_limites_bairros}. Sistema de referência SIRGAS 2000.}{mapa:mapa-violencia-territorial-homicidios-jovens-negros-ra-2024}

Em 2024, a taxa de homicídios de jovens negros apresentou diferenças importantes entre as Regiões Administrativas do Rio de Janeiro. Os maiores valores foram registrados em Madureira (9,18 por 100 mil habitantes), Pavuna (6,77), Cidade de Deus (6,54) e Centro (6,11). Em contraste, Tijuca e Santa Teresa apresentaram taxa igual a zero.

O mapa evidencia uma distribuição territorial desigual desse indicador, com maiores taxas concentradas principalmente em algumas áreas das Zonas Norte e Oeste. Assim como nos mapas anteriores, o dado deve ser interpretado como um indicador do contexto de violência territorial, e não como medida específica da violência sofrida por crianças na primeira infância.

\vertabelas{Tabela~\ref{tab:tabela-mapa-violencia-territorial-ra-2024}; ver Apêndice~\ref{ap:eixo-5}}

\section{Síntese do eixo}

\begin{sintese}{Direito ao Brincar}A violência territorial é muito desigual entre as regiões da cidade e atinge com mais força áreas da Zona Norte e a Cidade de Deus. Para as próximas edições, o eixo precisa de dados próprios da primeira infância, como a oferta e o estado de praças, parques e espaços de brincar, e de uma série histórica da violência territorial.\end{sintese}

\eixo{6}{7}{apple}
